\begin{abstract}
We present a controlled systems study of embodied-agent execution in Isaac Sim
that isolates two design axes at the planner/executor boundary:
planner-to-executor contract design and lightweight runtime validation. Across
navigation and tabletop manipulation tasks, we compare an under-specified
direct-action contract (P0), a typed tool-call contract (P1), and the same
typed contract with a brief self-check (P2), under either bare dispatch
(R0) or validation with a single retry (R1). In the fixed-runtime
action-interface ablation, P0/R0 fails on all four evaluated tasks, while
P1/R0 and P2/R0 succeed on all four and eliminate invalid actions. In the
fixed-contract runtime-only ablation, adding validation and one retry raises
success from 2/4 to 4/4 by converting invalid-first-action runs into
successful recoveries. Across the broader evaluation slices, the same ordering
remains stable: P0/R0 is weakest, P0/R1 recovers, P1 and P2 remain
strong, and P2 is associated with lower planner/tool overhead than P1 in
the covered comparisons. The contribution is controlled empirical evidence
that explicit action interfaces and lightweight runtime checks are meaningful
systems levers for embodied-agent execution in the tested Isaac Sim setting.
\end{abstract}
